% ============================================================================
%  PHYSICS NOTE: Celestial Holography and Calabi--Yau Quantum Groups
%
%  Status: research note / programme outline (not compiled into main.tex)
%  Dependencies: Vol I shadow obstruction tower (thqg_soft_graviton_theorems.tex),
%    Vol III quantum chiral algebras (quantum_chiral_algebras.tex),
%    toric CY3 / affine Yangians (toric_cy3_coha.tex),
%    E_2-chiral algebras (e2_chiral_algebras.tex)
% ============================================================================

\documentclass[11pt]{article}
\usepackage[T1]{fontenc}
\usepackage[utf8]{inputenc}
\usepackage{amsmath,amssymb,amsthm}
\usepackage{mathtools}
\usepackage{enumitem}
\usepackage{booktabs}
\usepackage[margin=1.2in]{geometry}
\usepackage{hyperref}

% Theorem environments
\newtheorem{theorem}{Theorem}[section]
\newtheorem{proposition}[theorem]{Proposition}
\newtheorem{lemma}[theorem]{Lemma}
\newtheorem{corollary}[theorem]{Corollary}
\newtheorem{conjecture}[theorem]{Conjecture}
\theoremstyle{definition}
\newtheorem{definition}[theorem]{Definition}
\newtheorem{example}[theorem]{Example}
\newtheorem{construction}[theorem]{Construction}
\theoremstyle{remark}
\newtheorem{remark}[theorem]{Remark}
\newtheorem{warning}[theorem]{Warning}

% Macros
\newcommand{\cA}{\mathcal{A}}
\newcommand{\cC}{\mathcal{C}}
\newcommand{\cH}{\mathcal{H}}
\newcommand{\cM}{\mathcal{M}}
\newcommand{\cO}{\mathcal{O}}
\newcommand{\cS}{\mathcal{S}}
\newcommand{\cV}{\mathcal{V}}
\newcommand{\cW}{\mathcal{W}}
\newcommand{\cZ}{\mathcal{Z}}
\newcommand{\frakg}{\mathfrak{g}}
\newcommand{\frakS}{\mathfrak{S}}
\newcommand{\HH}{\mathrm{HH}}
\newcommand{\Rep}{\mathrm{Rep}}
\newcommand{\Res}{\mathrm{Res}}
\newcommand{\Sh}{\mathrm{Sh}}
\newcommand{\Tr}{\mathrm{Tr}}
\newcommand{\Fuk}{\mathrm{Fuk}}
\newcommand{\Coh}{\mathrm{Coh}}
\newcommand{\CoHA}{\mathrm{CoHA}}
\newcommand{\DT}{\mathrm{DT}}
\newcommand{\CY}{\mathrm{CY}}
\newcommand{\Ran}{\mathrm{Ran}}
\newcommand{\End}{\mathrm{End}}
\newcommand{\Hom}{\mathrm{Hom}}
\newcommand{\id}{\mathrm{id}}
\newcommand{\op}{\mathrm{op}}
\newcommand{\MC}{\mathrm{MC}}

\title{Celestial Holography and Calabi--Yau Quantum Groups}
\author{Raeez Lorgat}
\date{April 2026 --- Research Note}

\begin{document}
\maketitle

\begin{abstract}
We connect celestial holography with the quantum vertex chiral groups of CY3s.  The Mellin transform of 4d scattering amplitudes produces a 2d celestial chiral algebra on $S^2$.  For 4d $\mathcal{N}=2$ theories from CY3 compactification, this algebra is the genus-0 sector of $G(X)$.  The simplest case, $\cW_{1+\infty} = Y(\widehat{\mathfrak{gl}}_1) = G(\mathbb{C}^3)$, recovers the $\mathfrak{w}_{1+\infty}$ symmetry of pure gravity.  The celestial OPE is the collision $r$-matrix $r(z) = \Res^{\mathrm{coll}}_{0,2}(\Theta_A)$; higher soft theorems correspond to higher-arity shadows of $\Theta_A$; higher-genus corrections control loop-level celestial amplitudes.
\end{abstract}

\tableofcontents

% ====================================================================
\section{Celestial amplitudes and the celestial chiral algebra}
\label{sec:celestial-amplitudes}
% ====================================================================

\subsection{The Mellin transform to the celestial sphere}

Four-dimensional scattering amplitudes in Minkowski space
$\mathbb{R}^{1,3}$ can be reformulated as correlation functions on
the celestial sphere $S^2 \cong \mathbb{P}^1$ at null infinity.
The passage is via the Mellin transform in the energy variable.

A massless particle of momentum $p^\mu$ can be parametrized as
\begin{equation}\label{eq:null-momentum-parametrization}
p^\mu = \varepsilon\, \omega\, q^\mu(z, \bar{z}),
\qquad
q^\mu(z,\bar{z}) = (1+z\bar{z},\; z+\bar{z},\;
-i(z-\bar{z}),\; 1-z\bar{z}),
\end{equation}
where $\omega > 0$ is the energy, $\varepsilon = \pm 1$ labels
incoming/outgoing, and $(z,\bar{z}) \in \mathbb{P}^1$ is the
point on the celestial sphere determined by the direction of
propagation.  An $n$-particle scattering amplitude
$\cA_n(\omega_i, z_i, \bar{z}_i)$ is Mellin-transformed to a
\emph{celestial amplitude}
\begin{equation}\label{eq:mellin-transform}
\widetilde{\cA}_n(\Delta_i, z_i, \bar{z}_i)
\;=\;
\prod_{i=1}^{n}
\int_0^\infty d\omega_i\; \omega_i^{\Delta_i - 1}\;
\cA_n(\omega_i, z_i, \bar{z}_i),
\end{equation}
where $\Delta_i \in \mathbb{C}$ is the \emph{conformal dimension}
of the $i$-th operator on the celestial sphere.  The Lorentz group
$\mathrm{SL}(2,\mathbb{C})$ acts on $\mathbb{P}^1$ by M\"obius
transformations, and $\widetilde{\cA}_n$ transforms as a
correlation function of primary operators of weights
$(\Delta_i, \bar{\Delta}_i)$ in a 2d CFT on $\mathbb{P}^1$.

\begin{definition}[Celestial chiral algebra]
\label{def:celestial-chiral-algebra}
The \emph{celestial chiral algebra} $\cA_{\mathrm{cel}}$ of a
4d massless theory is the holomorphic sector of the 2d CFT on
$\mathbb{P}^1$ whose correlation functions are the celestial
amplitudes~\eqref{eq:mellin-transform}.  This is a vertex algebra
on $\mathbb{P}^1$ whose OPE encodes the collinear singularities
of the 4d S-matrix.
\end{definition}

\begin{remark}[Collinear limits and the celestial OPE]
\label{rem:collinear-ope}
In the momentum-space S-matrix, collinear singularities arise when
two momenta become parallel: $p_i \parallel p_j$.  On the celestial
sphere, this corresponds to $z_i \to z_j$, which is exactly the
OPE limit.  The celestial OPE therefore encodes the universal
collinear behavior of the S-matrix:
\[
\cO_{\Delta_1}(z_1) \cdot \cO_{\Delta_2}(z_2)
\;\sim\;
\sum_k \frac{C_k(\Delta_1, \Delta_2)}{(z_1 - z_2)^{h_k}}
\;\cO_{\Delta_k}(z_2),
\]
where the OPE coefficients $C_k$ are determined by the collinear
splitting functions of the 4d theory.
\end{remark}

\subsection{Soft theorems as Ward identities}

The structural insight of celestial holography (Strominger 2014, He--Lysov--Mitra--Strominger 2015) is that soft theorems are Ward identities of the celestial chiral algebra.

\begin{itemize}
\item \textbf{Weinberg's leading soft graviton theorem} (1965):
  In the limit where one graviton's energy $\omega \to 0$, the
  amplitude factorizes universally.  On the celestial sphere,
  this becomes the Ward identity of a spin-2 current, which
  generates the supertranslation subalgebra of the BMS group.

\item \textbf{Cachazo--Strominger subleading soft theorem} (2014):
  The $\mathcal{O}(\omega^0)$ correction in the soft limit gives
  a subleading Ward identity, generating the superrotation
  (Virasoro) subalgebra of the extended BMS group.

\item \textbf{Higher soft theorems}: Each sub$^k$-leading soft
  theorem at $\mathcal{O}(\omega^{k-1})$ produces a Ward identity
  from a spin-$(k+1)$ current in the celestial chiral algebra.
\end{itemize}

\begin{proposition}[Soft theorems and the arity filtration]
\label{prop:soft-arity}
The soft graviton theorems are in bijection with the arity
filtration of the shadow obstruction tower of Vol~I.  The
sub$^k$-leading soft theorem is the Ward identity of the
arity-$(k+2)$ shadow $\Theta^{(k+2)}_\cA$, evaluated on
genus-0 conformal blocks via the shadow representation:
\[
\Sh_{0,n}(\Theta^{(k+2)}_\cA) \cdot \Phi = 0
\qquad
\text{(arity-$(k+2)$ soft Ward identity)}.
\]
\end{proposition}

This follows from the arity decomposition of the modular shadow
connection (Vol~I, Theorem~\ref*{thm:thqg-VI-flatness-by-arity}).
The leading soft theorem is $\Sh_{0,n}(\kappa(\cA))$, the
subleading is $\Sh_{0,n}(\mathfrak{C}(\cA))$, the sub-subleading
is $\Sh_{0,n}(\mathfrak{Q}(\cA))$, etc.

% ====================================================================
\section{The $\mathfrak{w}_{1+\infty}$ symmetry and $G(\mathbb{C}^3)$}
\label{sec:w-infinity-c3}
% ====================================================================

\subsection{Strominger's $\mathfrak{w}_{1+\infty}$}

Guevara, Himwich, Pate, and Strominger (2021) established that the
leading celestial symmetry algebra for Einstein gravity in 4d is
$\mathfrak{w}_{1+\infty}$, the area-preserving diffeomorphisms of
the celestial sphere --- the classical (large-$c$) limit of the
$\cW_{1+\infty}$ algebra.  At the quantum level, the celestial
symmetry algebra of gravity coupled to matter at central
charge $c$ is the full $\cW_{1+\infty}$ algebra.

The $\cW_{1+\infty}$ algebra at central charge $c$ is generated
by a tower of higher-spin currents $W^{(s)}(z)$ for $s = 1, 2,
3, \ldots$, with:
\begin{itemize}
\item $W^{(1)}(z)$: a $\mathrm{U}(1)$ current (spin 1),
\item $W^{(2)}(z) = T(z)$: the stress tensor (spin 2),
\item $W^{(s)}(z)$: spin-$s$ currents for all $s \geq 3$.
\end{itemize}
The OPE among these currents is determined by a single parameter
$c$ (the central charge) up to isomorphism, and the algebra is
freely generated in the sense that there are no null relations
among the generators at generic $c$.

\subsection{The identification $\cW_{1+\infty} = Y(\widehat{\mathfrak{gl}}_1) = G(\mathbb{C}^3)$}

Three seemingly independent constructions converge on the same
algebraic object:

\begin{theorem}[The triple identification]
\label{thm:triple-identification}
There are canonical isomorphisms
\begin{equation}\label{eq:triple-id}
\cW_{1+\infty}
\;\simeq\;
Y(\widehat{\mathfrak{gl}}_1)
\;\simeq\;
G(\mathbb{C}^3),
\end{equation}
where:
\begin{enumerate}[label=(\roman*)]
\item $\cW_{1+\infty}$ is the unique one-parameter family of vertex
  algebras freely generated by one current in each spin $s \geq 1$;
\item $Y(\widehat{\mathfrak{gl}}_1)$ is the affine Yangian of
  $\mathfrak{gl}_1$;
\item $G(\mathbb{C}^3)$ is the quantum vertex chiral group of the
  toric CY3 $\mathbb{C}^3$ (Definition~\ref{conj:qvcg}).
\end{enumerate}
\end{theorem}

\begin{proof}[Proof sketch]
The isomorphism $\cW_{1+\infty} \simeq Y(\widehat{\mathfrak{gl}}_1)$
is the theorem of Maulik--Okounkov (2019) and
Schiffmann--Vasserot (2013), established through the action of the
affine Yangian on the equivariant cohomology of the Hilbert scheme
of points on $\mathbb{C}^2$.

The isomorphism $Y(\widehat{\mathfrak{gl}}_1) \simeq G(\mathbb{C}^3)$
follows from the Schiffmann--Vasserot theorem (Vol~III,
Theorem~\ref*{thm:sv-c3}): the critical CoHA of $\mathbb{C}^3$
(Jordan quiver, cubic potential $W = \mathrm{tr}(XYZ - XZY)$)
is isomorphic to the positive half
$Y^+(\widehat{\mathfrak{gl}}_1)$, and the quantum vertex chiral
group $G(\mathbb{C}^3)$ is by definition the full affine Yangian
associated to the root datum of $\mathbb{C}^3$.
\end{proof}

\begin{remark}[Physical content]
\label{rem:physical-content-c3}
The identification~\eqref{eq:triple-id} has the following
physical reading:
\begin{enumerate}[label=(\alph*)]
\item The celestial symmetry algebra of 4d gravity
  ($\cW_{1+\infty}$) \emph{is} the quantum vertex chiral group
  of flat space ($G(\mathbb{C}^3)$).
\item The infinite tower of higher-spin celestial symmetries
  ($W^{(s)}$ for $s = 1, 2, 3, \ldots$) is the tower of generators
  of the affine Yangian, which in turn is the tower of BPS states
  in the DT theory of $\mathbb{C}^3$.
\item The celestial central charge $c$ is \emph{not} the modular
  characteristic $\kappa(G(\mathbb{C}^3))$ in general: $\kappa$ depends
  on the subalgebra (Heisenberg: $\kappa = k$, which equals $c$ at rank $d$, level $1$; Virasoro: $\kappa = c/2$;
  full $\cW_{1+\infty}$: divergent).  The DT genus-1 amplitude is
  controlled by the Heisenberg $\kappa = k$; see the anomaly cancellation
  note for the distinction.
\end{enumerate}
The simplest celestial symmetry algebra IS the quantum vertex
chiral group of the simplest CY3.
\end{remark}

% ====================================================================
\section{CY3 compactification and celestial chiral algebras}
\label{sec:cy3-celestial}
% ====================================================================

\subsection{4d $\mathcal{N}=2$ theories from CY3 compactification}

Type IIB string theory compactified on a Calabi--Yau threefold $X$
produces a 4d $\mathcal{N}=2$ supergravity theory whose low-energy
effective action is determined by the geometry of $X$.  The vector
multiplet moduli space is the complex structure moduli space
$\cM_{\mathrm{cs}}(X)$; the hypermultiplet moduli space involves
the complexified K\"ahler moduli together with RR flux data.

The BPS spectrum of the 4d theory is controlled by
Donaldson--Thomas invariants:
\begin{equation}\label{eq:bps-dt}
\Omega(\gamma; X) = \DT_\gamma(X),
\qquad \gamma \in H_3(X, \mathbb{Z}),
\end{equation}
where $\Omega(\gamma; X)$ is the BPS index (second helicity
supertrace) of charge $\gamma$, and $\DT_\gamma(X)$ is the
generalized DT invariant counting coherent sheaves / special
Lagrangians / D-branes of class $\gamma$.

\begin{conjecture}[CY3 celestial chiral algebra]
\label{conj:cy3-celestial}
For a 4d $\mathcal{N}=2$ theory $T_X$ arising from Type IIB on a
CY3 $X$, the celestial chiral algebra is
\begin{equation}\label{eq:cy3-celestial-identification}
\cA_{\mathrm{cel}}(T_X)
\;\simeq\;
G(X)^{g=0},
\end{equation}
the genus-0 (tree-level) sector of the quantum vertex chiral
group $G(X)$, with higher-genus corrections from the shadow obstruction tower.
\end{conjecture}

\begin{definition}[Quantum vertex chiral group]
The \emph{quantum vertex chiral group} $G(X)$ of a CY3 $X$ is the
$E_2$-chiral algebra whose:
\begin{enumerate}[label=(\roman*)]
\item \textbf{Root datum} is the generalized root datum
  $\mathcal{R}(X) = (\Lambda(X), \Delta^{\mathrm{re}},
  \Delta^{\mathrm{im}}_0, \Delta^{\mathrm{im}}_1, W(X), \rho,
  \mathrm{mult})$ extracted from the CY geometry (lattice from
  $K_0(\Coh(X))$, root multiplicities from DT invariants);
\item \textbf{Positive half} is the critical CoHA:
  $U(\mathfrak{n}_+) = \CoHA(X) = Y^+(\widehat{\frakg}_{Q_X})$
  for toric $X$;
\item \textbf{$E_2$-structure} is the braided monoidal structure on
  $\Rep^{E_2}(G(X))$, with braiding from the quantum group
  $R$-matrix of the affine super Yangian;
\item \textbf{Automorphic form} is the denominator identity of the
  associated BKM superalgebra: $\Phi_X$ (the DT partition function
  for toric CY3, the Igusa cusp form $\Delta_5$ for $K3 \times E$).
\end{enumerate}
\end{definition}

\subsection{The genus filtration and loop corrections}

The quantum vertex chiral group $G(X)$ has a natural genus
filtration inherited from the shadow obstruction tower of the
associated chiral algebra $A_X = \Phi(\Coh(X))$:
\begin{equation}\label{eq:genus-filtration}
G(X)^{g=0} \;\subset\; G(X)^{g \leq 1} \;\subset\;
G(X)^{g \leq 2} \;\subset\; \cdots \;\subset\; G(X).
\end{equation}

\begin{itemize}
\item \textbf{Genus 0} ($G(X)^{g=0}$): The \emph{tree-level}
  quantum vertex chiral group.  This is the sector controlled by
  the collision $r$-matrix and the genus-0 shadow obstruction tower.  For toric
  CY3, it is the Yangian $R$-matrix sector.  This is the celestial
  chiral algebra proper.

\item \textbf{Genus 1} ($G(X)^{g \leq 1}$): One-loop corrections.
  The genus-1 obstruction $\mathrm{obs}_1 = \kappa(A_X) \cdot
  \lambda_1$ produces the first quantum correction to the celestial
  OPE.  The modular characteristic $\kappa(A_X)$ controls the leading
  soft graviton factor (not to be confused with the celestial central
  charge $c$, which is a distinct quantity; see Remark~\ref{rem:physical-content-c3}).

\item \textbf{Genus $g$}: The genus-$g$ shadow obstruction tower component
  $\mathrm{obs}_g(A_X) = \kappa(A_X) \cdot \lambda_g$ (on the
  uniform-weight lane) produces $g$-loop corrections.  The
  full tower reproduces the all-loop celestial amplitude.
\end{itemize}

\begin{proposition}[Genus-0 celestial amplitudes from $G(X)$]
\label{prop:genus-0-celestial}
The tree-level celestial $n$-point amplitude of the 4d
$\mathcal{N}=2$ theory $T_X$ is the genus-0, $n$-point conformal
block of the quantum vertex chiral group:
\[
\widetilde{\cA}_n^{\mathrm{tree}}(\Delta_i, z_i)
\;=\;
\langle \cO_{\Delta_1}(z_1) \cdots \cO_{\Delta_n}(z_n)
\rangle_{G(X)^{g=0}}.
\]
The collinear singularity structure is controlled by the
collision $r$-matrix $r(z)$ of $G(X)$.
\end{proposition}

% ====================================================================
\section{The soft sector and the modular characteristic $\kappa$}
\label{sec:soft-sector-kappa}
% ====================================================================

\subsection{Leading soft theorems from $\kappa$}

In the Vol~I framework, the modular characteristic $\kappa(\cA)$
is the arity-2 projection of the universal MC element $\Theta_\cA$.
In the celestial context, $\kappa$ controls the leading soft
behavior.

% AP40: G(X) is conjectural for general CY3 (Conjecture CY-A_3);
% the claim below depends on its existence. Use conjecture environment.
\begin{conjecture}[Leading soft theorem = arity-2 projection]
\label{conj:leading-soft-kappa}
The Weinberg leading soft graviton theorem for the 4d theory $T_X$
is the arity-2 Ward identity of the shadow connection:
\begin{equation}\label{eq:leading-soft-kappa}
\Sh_{0,n}\bigl(\kappa(A_X)\bigr) \cdot \Phi
\;=\;
\sum_{i < j}
\frac{\kappa_{ij}}{z_i - z_j}\;\Phi
\;=\; 0,
\end{equation}
where $\kappa_{ij} = \kappa(A_X)(e_i, e_j)$ is the bilinear form
on charge vectors determined by the modular characteristic, and
$\Phi$ is a celestial correlator (horizontal section of
$\nabla^{\mathrm{hol}}_{0,n}$).

For $X = \mathbb{C}^3$: $\kappa(G(\mathbb{C}^3)) = \kappa(H_1) = 1$ (the
Heisenberg level, not $c/2 = 1/2$; see AP48), recovering the standard graviton soft
factor.
\end{conjecture}

\subsection{Higher soft theorems from higher-arity shadows}

Each arity level $r \geq 3$ of the shadow obstruction tower contributes a
higher soft theorem.

% AP40: Depends on conjectural G(X) for general CY3.
\begin{conjecture}[Higher soft theorems = higher shadows]
\label{conj:higher-soft-shadows}
The sub$^{r-2}$-leading soft theorem of the 4d theory $T_X$ is the
arity-$r$ Ward identity:
\begin{equation}\label{eq:higher-soft-shadow}
\Sh_{0,n}\bigl(\Theta_{A_X}^{(r)}\bigr) \cdot \Phi
\;=\;
-\sum_{\substack{s+t=r \\ s,t \geq 2}}
\bigl[\Sh_{0,n}(\Theta^{(s)}),\,
\Sh_{0,n}(\Theta^{(t)})\bigr] \cdot \Phi,
\end{equation}
where the right side involves the self-interaction of lower-arity
shadows.  The dictionary is:
\begin{center}
\small
\renewcommand{\arraystretch}{1.3}
\begin{tabular}{llll}
\toprule
\textit{Arity} & \textit{Shadow} &
\textit{Soft theorem} & \textit{Celestial current} \\
\midrule
$r=2$ & $\kappa(A_X)$ &
Weinberg leading & Supertranslation \\
$r=3$ & $\mathfrak{C}(A_X)$ &
Cachazo--Strominger subleading &
Superrotation (Virasoro) \\
$r=4$ & $\mathfrak{Q}(A_X)$ &
Sub-subleading &
Spin-3 $\cW$-current \\
$r$ & $\Theta_{A_X}^{(r)}$ &
Sub$^{r-2}$-leading &
Spin-$(r-1)$ $\cW$-current \\
\bottomrule
\end{tabular}
\end{center}
For $X = \mathbb{C}^3$ with shadow depth $r_{\max} = \infty$, the
full tower of soft theorems generates the $\cW_{1+\infty}$ algebra,
recovering the Strominger et al.\ result.
\end{conjecture}

\begin{remark}[Shadow depth and the celestial spectrum]
\label{rem:shadow-depth-celestial}
The shadow depth $r_{\max}(G(X))$ determines the celestial spin
content.  For CY3 geometries with finite shadow depth (if any
exist), only finitely many higher-spin currents appear; the
celestial chiral algebra is a truncation of $\cW_{1+\infty}$.
For generic CY3, we expect $r_{\max} = \infty$ (class
$\mathbf{M}$), giving the full infinite tower.  The physical
reason: the BPS spectrum of a generic CY3 is infinite, and each
BPS state contributes a celestial operator.
\end{remark}

% ====================================================================
\section{The celestial OPE as collision $r$-matrix}
\label{sec:celestial-ope-r-matrix}
% ====================================================================

\subsection{The collision $r$-matrix from $\Theta_A$}

The collision $r$-matrix of a quantum chiral algebra $A$ is the
arity-2 genus-0 projection of the universal MC element:
\begin{equation}\label{eq:r-matrix-definition}
r(z)
\;=\;
\Res^{\mathrm{coll}}_{0,2}(\Theta_A)
\;\in\;
A \otimes A \otimes k((z)),
\end{equation}
the residue of $\Theta_A$ along the collision divisor in
$\overline{\cM}_{0,3}$ (the moduli space with two coalescing points
and one output).  This is the $E_2$-chiral $R$-matrix of
Vol~III, Construction~\ref*{constr:cy-r-matrix}.

\begin{warning}[Bar kernel pole shift]
\label{warn:pole-shift}
The bar complex extracts residues along $d\log(z_i - z_j)$, not
along $dz/(z-w)$.  The $d\log$ measure absorbs one power of
$(z-w)$; hence $r(z)$ has pole orders one less than the OPE
(Vol~I, AP19).  For a Virasoro OPE with poles at $z^{-4}$,
$z^{-2}$, $z^{-1}$, the $r$-matrix has poles at $z^{-3}$, $z^{-1}$.
\end{warning}

\subsection{The celestial OPE}

% AP40: Depends on conjectural G(X) for general CY3.
\begin{conjecture}[Celestial OPE = collision $r$-matrix]
\label{conj:celestial-ope-r-matrix}
The celestial OPE of the 4d theory $T_X$ is determined by the
collision $r$-matrix of the quantum vertex chiral group $G(X)$:
\begin{equation}\label{eq:celestial-ope}
\cO_{\Delta_1}(z_1) \cdot \cO_{\Delta_2}(z_2)
\;\sim\;
r_{G(X)}(z_1 - z_2) \cdot
(\cO_{\Delta_1} \otimes \cO_{\Delta_2})(z_2),
\end{equation}
where $r_{G(X)}(z)$ is defined by~\eqref{eq:r-matrix-definition}
with $A = A_X$ the chiral algebra underlying $G(X)$.

The pole structure of the celestial OPE encodes the shadow depth:
\begin{enumerate}[label=(\roman*)]
\item The number of singular terms is $r_{\max}(G(X)) - 1$.
\item The leading \emph{simple} pole is determined by $\kappa(A_X)$:
  the $1/z$ coefficient of $r(z)$ is proportional to $\kappa(A_X)$.
  % AP19: For algebras with higher-order OPE poles (e.g.\ Virasoro),
  % r(z) has HIGHER poles as well (z^{-3} for Vir), not just z^{-1}.
  % The d-log kernel absorbs one power, so OPE pole z^{-n} becomes
  % r-matrix pole z^{-(n-1)}.  The statement r(z) ~ kappa/z + ...
  % is correct only for quadratic OPE (Heisenberg/KM).  For Virasoro:
  % r(z) = (c/2)/z^3 + 2T/z (AP19).
\item The subleading poles are determined by the higher shadows:
  the $(z_1 - z_2)^{-(k+1)}$ pole coefficient is
  $C_k = \Res^{\mathrm{coll}}_{0,2}(\Theta_{A_X}^{(k+2)})$.
\end{enumerate}
\end{conjecture}

\begin{remark}[Argument sketch]
The celestial OPE is the short-distance expansion of two operator
insertions on the celestial sphere $\mathbb{P}^1$.  In the
shadow-connection framework, this is the monodromy of the flat
connection $\nabla^{\mathrm{hol}}_{0,n}$ around the diagonal
$z_1 = z_2$.  By Vol~I (Theorem~\ref*{thm:thqg-VI-soft-ope}),
this monodromy is controlled by $\Sh_{0,n}(\Theta_A)$, and the
leading singular part is the collision residue
$\Res^{\mathrm{coll}}_{0,2}(\Theta_A) = r(z)$.  Each arity-$r$
component of $\Theta_A$ contributes a pole of order $r-1$ in the
celestial OPE.
\end{remark}

\begin{remark}[The quantum Yang--Baxter equation]
\label{rem:qybe-celestial}
The $r$-matrix $r_{G(X)}(z)$ satisfies the quantum Yang--Baxter
equation:
\[
r_{12}(z_1 - z_2)\, r_{13}(z_1 - z_3)\, r_{23}(z_2 - z_3)
\;=\;
r_{23}(z_2 - z_3)\, r_{13}(z_1 - z_3)\, r_{12}(z_1 - z_2).
\]
In the celestial context, this is the \emph{associativity of the
celestial OPE}: the OPE of three operators is independent of the
order in which the OPE limits are taken.  The QYBE is a consequence
of the flatness $(\nabla^{\mathrm{hol}})^2 = 0$, which in turn is
the MC equation for $\Theta_A$.
\end{remark}

% ====================================================================
\section{The full dictionary: CY3 geometry to celestial holography}
\label{sec:full-dictionary}
% ====================================================================

The passage from CY3 geometry to celestial holography factors
through the quantum vertex chiral group as an intermediary:

\begin{center}
\begin{tabular}{p{3cm}p{4cm}p{4.5cm}}
\toprule
\textit{CY3 geometry} &
\textit{Quantum vertex chiral group} &
\textit{Celestial holography} \\
\midrule
$X$ (CY threefold) &
$G(X)$ (quantum vertex chiral group) &
$\cA_{\mathrm{cel}}(T_X)$ (celestial chiral algebra) \\[4pt]
$\Lambda(X) = K_0(\Coh(X))$ &
Root lattice of $G(X)$ &
Charge lattice of celestial operators \\[4pt]
$\DT_\gamma(X)$ (DT invariants) &
Root multiplicities of $G(X)$ &
BPS celestial operator spectrum \\[4pt]
$\CoHA(X)$ (critical CoHA) &
$U(\mathfrak{n}_+) = $ positive half &
Holomorphic collinear sector \\[4pt]
$r(z) = \Res^{\mathrm{coll}}_{0,2}(\Theta_{A_X})$ &
$E_2$-chiral $R$-matrix &
Celestial OPE \\[4pt]
$\kappa(A_X)$ (modular characteristic) &
Arity-2 shadow &
Leading soft graviton theorem \\[4pt]
$\mathfrak{C}(A_X)$ (cubic shadow) &
Arity-3 shadow &
Subleading soft theorem \\[4pt]
$\Theta_{A_X}^{(r)}$ (arity-$r$ shadow) &
Shadow obstruction tower &
Sub$^{r-2}$-leading soft theorem \\[4pt]
$\mathrm{obs}_g(A_X)$ (genus-$g$ obstruction) &
Genus-$g$ shadow obstruction tower &
$g$-loop celestial amplitude \\[4pt]
$\Delta_5$ / $Z^{\DT}(X)$ (automorphic form) &
Denominator identity &
Celestial partition function \\
\bottomrule
\end{tabular}
\end{center}

\subsection{The automorphic correction as the shadow obstruction tower}

The central identification of the Vol~III programme (VISION.md) is:
\begin{equation}\label{eq:automorphic-shadow}
\boxed{%
\text{Automorphic correction}
\;=\;
\text{Shadow obstruction tower}
\;=\;
\text{Higher soft theorems.}}
\end{equation}

In detail: the passage from the ``naive'' algebra $\frakg$ (defined
by the real root Gram matrix alone) to the full BKM superalgebra
$\frakg_X$ (with all imaginary roots and their multiplicities)
is the automorphic correction.  In the Vol~I language, this is the
passage from the tree-level $r$-matrix $r(z)$ (arity 2) to the
full MC element $\Theta_A$ (all arities).  In the celestial
language, it is the passage from the leading soft theorem to the
full tower of soft theorems:

\begin{itemize}
\item \textbf{Arity 2} ($\kappa$): The leading coefficient =
  the Weyl vector contribution.  Celestially: the Weinberg soft
  factor.
\item \textbf{Arity 3} ($\mathfrak{C}$): The first imaginary root
  corrections.  Celestially: the Cachazo--Strominger subleading
  soft theorem.
\item \textbf{Arity 4} ($\mathfrak{Q}$): Higher corrections.
  Celestially: the sub-subleading soft theorem with contact
  coefficient $Q^{\mathrm{contact}}$.
\item \textbf{Full tower} ($\Theta_A$): The complete automorphic
  correction, whose denominator identity is the automorphic form.
  Celestially: the infinite tower of soft theorems whose
  consistency (MC equation) is the associativity of the celestial
  OPE.
\end{itemize}

\subsection{The $K3 \times E$ case: Igusa cusp form and celestial holography}

For the CY3 family $X = (S \times E)/(\mathbb{Z}/N\mathbb{Z})$
(K3 surface $S$, elliptic curve $E$):

\begin{example}[$K3 \times E$ celestial holography]
\label{ex:k3-e-celestial}
The quantum vertex chiral group $G(K3 \times E)$ is the
generalized BKM superalgebra $\frakg_{\Delta_5}$ with:
\begin{enumerate}[label=(\alph*)]
\item Three real simple roots $\delta_1, \delta_2, \delta_3$ with
  Gram matrix
  $\left(\begin{smallmatrix}
  2 & -2 & -2 \\ -2 & 2 & -2 \\ -2 & -2 & 2
  \end{smallmatrix}\right)$.
\item Imaginary root multiplicities
  $\mathrm{mult}(\alpha) = f(nm, l)$ where $f(n,l)$ are Fourier
  coefficients of the weak Jacobi form $\phi_{0,1}$ (the K3
  elliptic genus).
\item Denominator identity = the Igusa cusp form $\Delta_5$ of
  weight~5 for $\mathrm{Sp}_4(\mathbb{Z})$.
\end{enumerate}

The celestial interpretation:
\begin{itemize}
\item The modular characteristic
  $\kappa(G(K3 \times E)) = 5$ (the weight of $\Delta_5$).
\item The celestial chiral algebra $\cA_{\mathrm{cel}}$ is the
  genus-0 sector of $\frakg_{\Delta_5}$.
\item The $\mathrm{Sp}_4(\mathbb{Z})$ symmetry of the Siegel
  modular form IS the $E_2$-braiding --- the braided monoidal
  structure on $\Rep^{E_2}(G(K3 \times E))$.
\item The shadow obstruction tower encodes the product side of the denominator
  identity:
  \[
  \frac{1}{64}\,\Delta_5
  \;=\;
  e^{-2\pi i(\rho, z)}
  \prod_{\alpha \in \Delta_+}
  \bigl(1 - e^{-2\pi i(\alpha, z)}\bigr)^{\mathrm{mult}(\alpha)}.
  \]
  Each factor in the product corresponds to a shadow obstruction tower
  component: the real roots ($\mathrm{mult} = 1$) contribute at
  low arity, and the imaginary roots contribute at higher arities
  with multiplicities from $\phi_{0,1}$.
\end{itemize}
\end{example}

\subsection{The toric CY3 case: affine super Yangians}

\begin{example}[Toric CY3 celestial holography]
\label{ex:toric-celestial}
For a toric CY3 $X$ without compact 4-cycles:
\begin{enumerate}[label=(\alph*)]
\item The quantum vertex chiral group
  $G(X) = Y(\widehat{\frakg}_{Q_X})$, the affine super Yangian
  of the toric quiver.
\item The celestial chiral algebra is the genus-0 sector: the
  Yangian $R$-matrix determines the celestial OPE.
\item For $X = \mathbb{C}^3$: $G(\mathbb{C}^3) = \cW_{1+\infty}$,
  recovering the Strominger et al.\ result directly.
\item For the resolved conifold: the toric diagram has two vertices,
  the quiver is Klebanov--Witten (two nodes, bidirectional arrows),
  and $G(X)$ is the affine super Yangian of the corresponding
  quiver Lie superalgebra.  The celestial chiral algebra has
  two families of higher-spin currents, one for each node.
\end{enumerate}
\end{example}

% ====================================================================
\section{Structural theorems}
\label{sec:structural}
% ====================================================================

\subsection{The celestial shadow connection}

The shadow connection of Vol~I, specialized to genus 0 and evaluated
on the quantum vertex chiral group $G(X)$, gives the celestial
shadow connection:

\begin{definition}[Celestial shadow connection]
\label{def:celestial-shadow-connection}
The \emph{celestial shadow connection} of the 4d theory $T_X$ is
\begin{equation}\label{eq:celestial-shadow-connection}
\nabla^{\mathrm{cel}}_{n}
\;:=\;
\nabla^{\mathrm{hol}}_{0,n}\bigl|_{A = A_X}
\;=\;
d \;-\; \Sh_{0,n}(\Theta_{A_X}),
\end{equation}
a flat connection on the space of $n$-point celestial amplitudes
$\cV_{0,n}(A_X)$.  Its flatness $(\nabla^{\mathrm{cel}}_n)^2 = 0$
is the MC equation for $\Theta_{A_X}$.
\end{definition}

\begin{theorem}[Celestial Ward identities from the shadow obstruction tower]
\label{thm:celestial-ward}
The horizontal sections of $\nabla^{\mathrm{cel}}_n$ are the
celestial correlators.  The arity decomposition of
$\nabla^{\mathrm{cel}}_n$ produces the full tower of soft
graviton theorems:
\begin{equation}\label{eq:celestial-ward-tower}
\nabla^{\mathrm{cel}}_n
\;=\;
\nabla^{(2)}_n + A^{(3)}_n + A^{(4)}_n + \cdots,
\end{equation}
where $\nabla^{(2)}_n$ is the leading (supertranslation)
connection and $A^{(r)}_n$ for $r \geq 3$ are the higher soft
perturbations.  The arity-$r$ flatness equation
$\nabla^{(2)}(A^{(r)}) + \sum_{s+t=r}[A^{(s)}, A^{(t)}] +
o_r = 0$ is the sub$^{r-2}$-leading soft theorem with all
nonlinear corrections from lower arities.
\end{theorem}

\subsection{The celestial Koszul dual}

\begin{conjecture}[Celestial Koszul duality]
\label{conj:celestial-koszul}
The Koszul dual $G(X)^!$ of the quantum vertex chiral group
controls the celestial amplitudes of the \emph{dual} 4d theory.
In the $E_2$-chiral Koszul duality framework (Vol~III,
Conjecture~\ref*{conj:e2-koszul}):
\begin{equation}\label{eq:celestial-koszul}
\Rep^{E_2}(G(X))
\;\simeq\;
\Rep^{E_2}(G(X)^!)^{\mathrm{rev}},
\end{equation}
where $\mathrm{rev}$ denotes the reverse braiding.

For $X = \mathbb{C}^3$: $\cW_{1+\infty}$ is self-dual (the
Feigin--Frenkel involution acts trivially at the self-dual level),
and the celestial chiral algebra of 4d gravity is self-dual.
\end{conjecture}

\subsection{The $E_2$-structure and braided celestial amplitudes}

\begin{proposition}[The $E_2$-braiding is the celestial crossing symmetry]
\label{prop:e2-crossing}
The $E_2$-structure on $G(X)$ (the braided monoidal structure on
$\Rep^{E_2}(G(X))$) is the crossing symmetry of celestial
amplitudes.  The braiding
\[
\sigma_{V,W} \colon V \otimes W \;\xrightarrow{\sim}\; W \otimes V
\]
in $\Rep^{E_2}(G(X))$ corresponds to the exchange
$z_i \leftrightarrow z_j$ of celestial operator positions,
implemented by the monodromy of $\nabla^{\mathrm{cel}}$ around the
exchange path in configuration space.
\end{proposition}

The $R$-matrix
$R(z) \in \End(V \otimes W)((z))$ satisfying the QYBE is the
algebraic encoding of this crossing symmetry.  For the Yangian
$Y(\widehat{\mathfrak{gl}}_1)$, the $R$-matrix is the rational
$R$-matrix $R(u) = 1 + P/u$ (where $P$ is the permutation
operator), giving the simplest crossing kernel for graviton
celestial amplitudes.

% ====================================================================
\section{Loop corrections and the genus tower}
\label{sec:loop-corrections}
% ====================================================================

\subsection{The genus expansion of celestial amplitudes}

The shadow obstruction tower of the quantum vertex chiral group produces a
genus expansion for celestial amplitudes that corresponds to the
loop expansion of the 4d S-matrix:

\begin{equation}\label{eq:genus-expansion-celestial}
\widetilde{\cA}_n(\Delta_i, z_i)
\;=\;
\sum_{g=0}^{\infty}
\hbar^{2g-2}\;
\widetilde{\cA}_n^{(g)}(\Delta_i, z_i),
\end{equation}
where $\widetilde{\cA}_n^{(g)}$ is the genus-$g$ celestial
amplitude, computed from the genus-$g$ shadow obstruction tower of $G(X)$.

\begin{theorem}[Loop corrections from the shadow obstruction tower]
\label{thm:loop-corrections}
The genus-$g$ celestial amplitude of the 4d theory $T_X$ is
controlled by the genus-$g$ obstruction class:
\begin{equation}\label{eq:genus-g-celestial}
\widetilde{\cA}_n^{(g)}
\;=\;
\int_{\overline{\cM}_{g,n}}
\mathrm{obs}_g(A_X) \cdot
\langle \cO_{\Delta_1}(z_1) \cdots \cO_{\Delta_n}(z_n)
\rangle_{g,n},
\end{equation}
where $\mathrm{obs}_g(A_X) = \kappa(A_X) \cdot \lambda_g$ on the
uniform-weight lane, and the integral is over the moduli space
$\overline{\cM}_{g,n}$ of stable genus-$g$ curves with $n$ marked
points.

The all-genera generating function is:
\begin{equation}\label{eq:all-genera-celestial}
F_{\mathrm{cel}}(\hbar)
\;=\;
\sum_{g=1}^{\infty} \hbar^{2g}\;
F_g(A_X)
\;=\;
\kappa(A_X) \cdot
\bigl(\hat{A}(i\hbar) - 1\bigr),
\end{equation}
where $\hat{A}$ is the $\hat{A}$-genus and $F_g(A_X) =
\kappa(A_X) \cdot \int_{\overline{\cM}_g} \lambda_g$ are the
shadow obstruction tower amplitudes of Vol~I.
\end{theorem}

\begin{remark}[The $\hbar$ convention]
\label{rem:hbar-convention}
The factor $\hbar^{2g-2}$ in~\eqref{eq:genus-expansion-celestial}
uses the Vol~I convention where $\hbar^{2g-2}$ weights the genus-$g$
contribution.  The Feynman-section convention $\hbar^{g-1}$ is
related by $\hbar_{\mathrm{there}} = \hbar_{\mathrm{here}}^2$.
The leading term $g=0$ at $\hbar^{-2}$ is the tree-level celestial
amplitude.
\end{remark}

\subsection{Interpretation: GW invariants and celestial loop amplitudes}

The shadow obstruction tower amplitudes $F_g(A_X)$ involve the Hodge class
$\lambda_g$ integrated over $\overline{\cM}_g$.  For a CY3 $X$,
the Gromov--Witten invariants $\mathrm{GW}_{g,n,\beta}(X)$ (counting
holomorphic curves of genus $g$ and class $\beta$ in $X$ through
$n$ constraints) are related to the shadow obstruction tower via the CY modular
characteristic (Vol~III, Theorem CY-D).  The celestial loop
amplitudes therefore encode the enumerative geometry of $X$:
higher-loop celestial amplitudes ``count curves'' in the CY3.

% ====================================================================
\section{Comparison with the celestial holography literature}
\label{sec:comparison}
% ====================================================================

\subsection{Relation to Costello--Paquette}

Costello and Paquette (2022) construct the celestial chiral algebra
of self-dual gravity as a specific vertex algebra (the ``BCOV
chiral algebra''), derived from the Kodaira--Spencer theory of
gravity on a CY3.  In our framework:

\begin{itemize}
\item The Kodaira--Spencer / BCOV theory on $X$ is the
  ``closed-string'' side of the CY category $\Coh(X)$.
\item The BCOV chiral algebra on $\mathbb{P}^1$ is a model for the
  genus-0 sector of $G(X)$.
\item The $\cW_{1+\infty}$ symmetry discovered by Costello--Paquette
  for $X = \mathbb{C}^3$ is recovered as
  $G(\mathbb{C}^3)^{g=0} = Y(\widehat{\mathfrak{gl}}_1) =
  \cW_{1+\infty}$.
\item The higher-genus corrections that are absent in the
  Costello--Paquette construction (which works at the self-dual /
  tree level) are supplied by the shadow obstruction tower.
\end{itemize}

\subsection{Relation to Guevara--Himwich--Pate--Strominger}

Guevara, Himwich, Pate, and Strominger identified the leading
celestial symmetry as $\mathfrak{w}_{1+\infty}$ (the classical
limit of $\cW_{1+\infty}$) and conjectured its promotion to the
full quantum $\cW_{1+\infty}$ at finite $c$.  Our framework:

\begin{itemize}
\item Provides the algebraic mechanism: the quantum vertex chiral
  group $G(\mathbb{C}^3) = \cW_{1+\infty}$ at all $c$ simultaneously
  (the affine Yangian is defined for arbitrary $c$).
\item Explains the origin of $\cW_{1+\infty}$: it arises as the
  quantum vertex chiral group of the simplest CY3, not as an
  \emph{ad hoc} symmetry algebra.
\item Generalizes to arbitrary CY3: for each $X$, there is a
  celestial chiral algebra $G(X)^{g=0}$ that extends
  $\cW_{1+\infty}$ according to the geometry of $X$.
\end{itemize}

\subsection{Relation to Bu--Kim--Li--Zhang}

Bu, Kim, Li, and Zhang (2024) relate the celestial OPE to the
Yangian $R$-matrix in the context of self-dual Yang--Mills theory.
In our framework, this is the statement that the celestial OPE
is the collision $r$-matrix:
$r(z) = \Res^{\mathrm{coll}}_{0,2}(\Theta_A)$.

% ====================================================================
\section{Conjectures and open problems}
\label{sec:conjectures}
% ====================================================================

\begin{conjecture}[CY3 celestial universality]
\label{conj:cy3-universality}
For every compact CY3 $X$, there exists a quantum vertex chiral
group $G(X)$ whose genus-0 sector is the celestial chiral algebra
of the 4d $\mathcal{N}=2$ theory $T_X$:
\[
\cA_{\mathrm{cel}}(T_X)
\;\simeq\;
G(X)^{g=0}.
\]
The construction applies to all CY3 categories (not just toric or
$K3 \times E$) via the CY-to-chiral functor $\Phi$ of Vol~III,
Theorem CY-A.
\end{conjecture}

\begin{conjecture}[Wall-crossing = gauge equivalence]
\label{conj:wall-crossing}
The Kontsevich--Soibelman wall-crossing formula for DT invariants,
applied to the root multiplicities of $G(X)$, is a gauge
equivalence of MC elements in the modular convolution algebra:
\[
\Theta_{A_X}|_{t_+}
\;\sim_{\mathrm{gauge}}\;
\Theta_{A_X}|_{t_-},
\]
where $t_\pm$ are stability parameters on opposite sides of a
wall of marginal stability.  The celestial chiral algebra changes
by a gauge transformation, not an isomorphism --- the celestial
OPE coefficients jump, but the underlying algebraic structure
(the QYBE, the MC equation) is preserved.
\end{conjecture}

\begin{conjecture}[Celestial shadow depth]
\label{conj:celestial-shadow-depth}
The shadow depth $r_{\max}(G(X))$ of the quantum vertex chiral
group of a CY3 $X$ is:
\begin{enumerate}[label=(\roman*)]
\item $r_{\max} = \infty$ for all smooth compact CY3s (the BPS
  spectrum is generically infinite);
\item $r_{\max} = \infty$ for all toric CY3s;
\item Finite shadow depth can occur only for degenerate limits
  (e.g.\ orbifold singularities where the BPS spectrum truncates).
\end{enumerate}
This implies that the celestial chiral algebra of a generic
CY3 compactification has the full $\cW_{1+\infty}$-type structure
--- an infinite tower of higher-spin symmetries.
\end{conjecture}

\begin{conjecture}[Eight celestial chiral algebras from $K3 \times E$]
\label{conj:eight-celestial}
The eight diagonal divisor modular forms of the Igusa cusp form
programme correspond to eight distinct quantum vertex chiral groups
$G(X_N)$ for $X_N = (S \times E)/(\mathbb{Z}/N\mathbb{Z})$,
$N = 1, \ldots, 8$, with:
\begin{enumerate}[label=(\roman*)]
\item Denominator identities = the eight dd-modular forms;
\item Root multiplicities from twisted-twined K3 elliptic genera;
\item Celestial chiral algebras $G(X_N)^{g=0}$ related by
  orbifolding of the $\mathbb{Z}/N\mathbb{Z}$ action.
\end{enumerate}
The Koszul duality relations among these eight algebras should
manifest as ``celestial S-dualities'' exchanging weak and strong
coupling in the 4d theories.
\end{conjecture}

\begin{conjecture}[Non-toric CY3 and the BPS root datum]
\label{conj:non-toric}
For a non-toric CY3 $X$ (e.g.\ the quintic threefold in
$\mathbb{P}^4$), the quantum vertex chiral group $G(X)$ exists
with root datum $\mathcal{R}(X)$ determined by:
\begin{enumerate}[label=(\roman*)]
\item Lattice: $\Lambda(X) = K_0(\Coh(X))$ (or
  $K_0(\Fuk(X))$ on the mirror side);
\item Root multiplicities: DT invariants $\DT_\gamma(X)$ (or
  GV / Gopakumar--Vafa invariants);
\item The ``BPS algebra'' of Joyce, Kontsevich--Soibelman, and
  Meinhardt--Reineke as the positive half;
\item Denominator identity: the DT generating function
  $Z^{\DT}(X, q)$.
\end{enumerate}
The key challenge is that without the toric combinatorics, the
quiver description is not available; the critical CoHA must be
replaced by a more general motivic Hall algebra construction.
\end{conjecture}

% ====================================================================
\section{Summary: the CY3 celestial programme}
\label{sec:summary}
% ====================================================================

The programme developed in this note connects three circles of
ideas through the quantum vertex chiral group:

\begin{enumerate}[label=(\arabic*)]
\item \textbf{CY3 geometry} (DT invariants, BPS spectra, K3
  elliptic genera, automorphic forms) provides the root datum
  $\mathcal{R}(X)$ of the quantum vertex chiral group $G(X)$.

\item \textbf{The shadow obstruction tower / bar-cobar machine} (Vol~I)
  provides the algebraic engine: the universal MC element
  $\Theta_A$, its arity filtration (the shadow obstruction tower),
  and the flat shadow connection $\nabla^{\mathrm{hol}}$.

\item \textbf{Celestial holography} provides the physical
  interpretation: 4d scattering amplitudes on the celestial
  sphere, organized by soft theorems that are the Ward identities
  of the shadow connection.
\end{enumerate}

The central results and identifications are:

\begin{enumerate}[label=(\alph*)]
\item The simplest celestial symmetry algebra, $\cW_{1+\infty}$,
  IS the quantum vertex chiral group of $\mathbb{C}^3$:
  $\cW_{1+\infty} = Y(\widehat{\mathfrak{gl}}_1) = G(\mathbb{C}^3)$.

\item For CY3 compactifications, the celestial chiral algebra is
  the genus-0 sector of $G(X)$, with higher-genus corrections
  from the shadow obstruction tower controlling loop amplitudes.

\item The celestial OPE is the collision $r$-matrix:
  $r(z) = \Res^{\mathrm{coll}}_{0,2}(\Theta_{A_X})$.

\item The leading soft theorem is the arity-2 projection $\kappa$.
  Higher soft theorems are higher-arity shadows.

\item The automorphic correction (passage from real roots to the
  full BKM with imaginary roots) IS the shadow obstruction tower
  (passage from arity-2 $\kappa$ to the full $\Theta_A$)
  IS the tower of soft theorems.

\item The $E_2$-braiding of $G(X)$ is the crossing symmetry of
  celestial amplitudes, with the quantum Yang--Baxter equation as
  the algebraic encoding of OPE associativity.
\end{enumerate}

\end{document}
